% uid: ""
% school: "lycee"
% level: "1ere"
% chapter: "Activités numériques"
% title: "Écriture scientifique"
% difficulty: "Moyen"
% tags: ["ecriture-scientifique", "puissances-de-10", "puissances", "nombre-de-chiffres"]
% provenance: "own"
% batch: "2026-11-22-revision-1ere-A"

% ---STATEMENT---
Écrire en écriture scientifique :
\begin{enumerate}
  \item $A=\dfrac{0{,}00045\times8\,000\,000\times(2{,}5\times10^{-3})^2}{0{,}0000125\times4\times10^{5}}$ ;
  \item $B=4^{15}\times25^{14}$ ;
  \item $C=2^{2026}\times5^{2020}$ ; combien de chiffres $C$ possède-t-il en écriture décimale ?
  \item \textbf{Application.} La Terre est à $1{,}5\times10^{8}$ km du Soleil et la lumière parcourt $3\times10^{5}$ km par seconde. Combien de temps la lumière du Soleil met-elle pour nous atteindre (en minutes et secondes) ?
  \item \textbf{Application.} Proxima du Centaure est à $4\times10^{13}$ km de nous. Combien d'années la lumière met-elle pour nous en parvenir (au dixième près) ? On prendra $1$ an $=3{,}15\times10^{7}$ s.
\end{enumerate}

% ---SOLUTION---
\begin{prop}[Écriture scientifique et puissances]
Écriture scientifique : \ib{$a\times10^{n}$ avec $1\leq a<10$ et $n\in\mathbb Z$}.

Règles : \ib{$10^{m}\times10^{n}=10^{m+n}$}, \ib{$(x^{m})^{n}=x^{mn}$}, \ib{$(xy)^{n}=x^{n}y^{n}$} et $\dfrac{10^{m}}{10^{n}}=10^{m-n}$.
\end{prop}
\begin{enumerate}
  \item On écrit chaque nombre en écriture scientifique :
  \[0{,}00045=4{,}5\times10^{-4}\qquad8\,000\,000=8\times10^{6}\qquad0{,}0000125=1{,}25\times10^{-5}\]
  \[(2{,}5\times10^{-3})^2=6{,}25\times10^{-6}\]
  \textbf{Numérateur :} on regroupe les décimaux et les puissances de $10$ :
  \[(4{,}5\times8\times6{,}25)\times10^{-4+6-6}=225\times10^{-4}\]
  \textbf{Dénominateur :}
  \[(1{,}25\times4)\times10^{-5+5}=5\]
  \textbf{Quotient :}
  \[A=\dfrac{225\times10^{-4}}{5}=45\times10^{-4}=\ib{$4{,}5\times10^{-3}$}\]

  \item On exprime tout avec des puissances de $2$ et de $5$ :
  \[4^{15}=(2^2)^{15}=2^{30}\qquad25^{14}=(5^2)^{14}=5^{28}\]
  On sépare $2^{30}=2^{2}\times2^{28}$ pour former des $10$ :
  \[B=2^{2}\times2^{28}\times5^{28}=4\times(2\times5)^{28}=\ib{$4\times10^{28}$}\]

  \item On sépare $2^{2026}=2^{6}\times2^{2020}$ :
  \[C=2^{6}\times2^{2020}\times5^{2020}=64\times10^{2020}\]
  Mais $64$ n'est pas compris entre $1$ et $10$ : on écrit $64=6{,}4\times10$.
  \[C=\ib{$6{,}4\times10^{2021}$}\]
  En écriture décimale, $C$ s'écrit $64$ suivi de $2020$ zéros.

  Le nombre de chiffres est $2+2020=\boxed{2022}$.

  \item On utilise la formule \og temps $=\dfrac{\text{distance}}{\text{vitesse}}$ \fg{} :
  \[t=\dfrac{1{,}5\times10^{8}}{3\times10^{5}}=0{,}5\times10^{3}=5\times10^{2}\ \text{s}=500\ \text{s}\]
  Comme $500=8\times60+20$ :
  \[t=\ib{$8\ \text{min}\ 20\ \text{s}$}\]

  \item Temps en secondes :
  \[t=\dfrac{4\times10^{13}}{3\times10^{5}}=\dfrac43\times10^{8}\approx1{,}33\times10^{8}\ \text{s}\]
  On convertit en années :
  \[\dfrac{1{,}33\times10^{8}}{3{,}15\times10^{7}}=\dfrac{13{,}3}{3{,}15}\approx4{,}2\]
  La lumière met environ \ib{$4{,}2$ ans}.
\end{enumerate}
\textbf{Idée :} pour des puissances de bases différentes, on cherche à former des $10=2\times5$.

\textbf{Piège :} laisser $64\times10^{2020}$ : ce n'est pas une écriture scientifique, et le nombre de chiffres est alors faux ($2022$ et non $2021$).
